% ============================================================
%  Monochromatic Union-Closed Families in the Boolean Lattice
% ============================================================
\documentclass[reqno]{aomart}

\usepackage{cmap}
\usepackage{amsmath,amsfonts,mathtools,amssymb,amsthm}
\usepackage[T1]{fontenc}
\usepackage{enumitem}
\hypersetup{colorlinks=true,linkcolor=blue,citecolor=red,urlcolor=blue}
\usepackage{booktabs,array}

\linespread{1.08}
\emergencystretch=2em
\raggedbottom
\setlength{\parskip}{3pt}

\numberwithin{equation}{section}

\theoremstyle{plain}
\newtheorem{theorem}{Theorem}[section]
\newtheorem{lemma}[theorem]{Lemma}
\newtheorem{proposition}[theorem]{Proposition}
\newtheorem{corollary}[theorem]{Corollary}
\newtheorem{question}[theorem]{Question}

\theoremstyle{definition}
\newtheorem{definition}[theorem]{Definition}

\theoremstyle{remark}
\newtheorem{remark}[theorem]{Remark}

\newcommand{\NN}{\mathbb{N}}
\newcommand{\RR}{\mathbb{R}}
\newcommand{\Ex}{\mathbb{E}}
\newcommand{\Prob}{\mathbb{P}}
\newcommand{\cF}{\mathcal{F}}
\newcommand{\cL}{\mathcal{L}}
\newcommand{\cY}{\mathcal{Y}}
\newcommand{\cJ}{\mathcal{J}}
\newcommand{\eps}{\varepsilon}
\DeclareMathOperator{\ind}{ind}
\DeclareMathOperator{\vc}{vc}
\DeclareMathOperator{\Down}{Down}
\DeclareMathOperator{\width}{w}

\title{Monochromatic Union-Closed Families in the Boolean Lattice}

\author[Deep Bhattacharjee]{Deep Bhattacharjee\textsuperscript{\dag}}
\address{Formerly, Electro-Gravitational Space Propulsion Laboratory (EGSPL),
  Bhubaneswar, Odisha 751030, India}
\email{d.bhattacharjee@erl-forschung.de}
\email{itsdeep@live.com}
\orcid{\href{https://orcid.org/0000-0003-0466-750X}{0000-0003-0466-750X}}
\thanks{\textsuperscript{\dag}\,Corresponding author.
\textsuperscript{\dag,\ddag}\,Deep Bhattacharjee and Priyabrata Mandal
contributed equally to this work and are joint first authors.}

\author[Priyabrata Mandal]{Priyabrata Mandal\textsuperscript{\ddag}}
\address{Department of Mathematics, Bioinformatics and Computer
  Applications, Maulana Azad National Institute of Technology Bhopal,
  Bhopal 462003, India}
\email{priyabrata@manit.ac.in}
\orcid{\href{https://orcid.org/0000-0001-6472-6239}{0000-0001-6472-6239}}

\author{Ushashi Bhattacharya}
\address{Formerly, National Taiwan University, Taipei 106319, Taiwan}
\email{rai.bhattacharya.in@gmail.com}
\orcid{\href{https://orcid.org/0009-0002-2254-3914}{0009-0002-2254-3914}}

\subjclass[2020]{Primary 05D05; Secondary 05D10, 06D05}
\keyword{union-closed families}
\keyword{Ramsey theory}
\keyword{Vapnik--Chervonenkis dimension}
\date{September 2026}

\copyrightnote{}
\fancypagestyle{firstpage}{%
  \fancyhf{}%
  \cfoot{\footnotesize\thepage}}

\begin{document}
\sloppy
\begin{abstract}
Let $f(n)$ be the largest integer such that every $2$-colouring of the
subsets of $\{1,\dots,n\}$ contains a monochromatic family of at least
$f(n)$ sets closed under union and intersection, and let $F(n)$ be the
analogous quantity for families closed under union only. Erd\H{o}s and
Ulam asked for the order of magnitude of both functions; Erd\H{o}s
recorded the trivial bound $f(n)\ge(n+1)/2$, wrote that there was no
plausible conjecture for the true order of $f(n)$, and conjectured that
$F(n)>n^{c}$ for every constant $c$ and that $F(n)<(1+\eps)^n$ for every
$\eps>0$. Random colourings give $f(n)\le n^2+n+1$ and
$f(n)<12n(\log_2(n+2))^2$, and, through the Sauer--Shelah lemma,
$F(n)\le\sum_{j<d}\binom nj$ whenever $2^d>nd+2$, so that
$F(n)\le n^{(1+o(1))\log_2n}$; this proves the second conjecture. The
limits of the random method are located: a random colouring contains, in
expectation, a monochromatic union-closed family of size
$n^{(1+o(1))\log_2\log_2n}$ and, with probability at least $1/2$, a
monochromatic chain of $n$ sets. Exact values are computed by a
counterexample-guided satisfiability search: $f(n)=\lceil(n+1)/2\rceil$ for
$1\le n\le10$, so that the trivial lower bound is sharp in this range, and
$F(n)$ is determined for $n\le7$; for $n\le10$ the extremal colourings for
$f$ can be chosen to depend only on the cardinality of the set, and these
colourings are listed. The upper bound for $F(n)$, the resolution of the
second conjecture, and the trivial lower bound are formally verified in
Lean~4 with Mathlib.
\end{abstract}

\maketitle

\section{Introduction}
\label{sec:intro}

Write $[n]=\{1,\dots,n\}$ and $2^{[n]}$ for the family of all subsets of
$[n]$. A \emph{$2$-colouring} is a map $\chi:2^{[n]}\to\{0,1\}$, and a
family $\cF\subseteq2^{[n]}$ is \emph{monochromatic} if $\chi$ is constant
on $\cF$. A family is \emph{union-closed} if $A\cup B\in\cF$ whenever
$A,B\in\cF$, and it is a \emph{sublattice} if it is closed under both union
and intersection. Every nonempty family with one member is closed under
both operations, so the following quantities are well defined:
\begin{align*}
f(n)&=\min_{\chi}\ \max\bigl\{|\cL|:\ \cL\subseteq2^{[n]}
\text{ a monochromatic sublattice}\bigr\},\\
F(n)&=\min_{\chi}\ \max\bigl\{|\cF|:\ \cF\subseteq2^{[n]}
\text{ a monochromatic union-closed family}\bigr\},
\end{align*}
the minima being over all $2$-colourings of $2^{[n]}$. Clearly
$F(n)\ge f(n)$. A maximal chain $\emptyset\subset\{1\}\subset\{1,2\}
\subset\dots\subset[n]$ has $n+1$ members, at least $\lceil(n+1)/2\rceil$
of one colour, and every subchain is a sublattice; hence
\begin{equation}
\label{eq:trivial}
F(n)\ \ge\ f(n)\ \ge\ \Bigl\lceil\frac{n+1}2\Bigr\rceil .
\end{equation}

The problem of estimating $f(n)$ and $F(n)$ was raised by Erd\H{o}s and
Ulam; Erd\H{o}s recorded it in \cite{Er78}, where he wrote that
\eqref{eq:trivial} was all they could say about $f(n)$ and that they had
no plausible conjecture for its true order of magnitude. For $F(n)$ he
conjectured that $F(n)>n^{c}$ for every constant $c$ and all $n>n_0(c)$,
and that $F(n)<(1+\eps)^n$ for every $\eps>0$ and all large $n$; he also
recorded, without reference, a result of Howorka that the first
conjecture holds for colourings in which all sets of the same cardinality
receive the same colour. The problem is Problem 1183 in Bloom's database
\cite{Bl26}, where it is listed as open with no further progress.

\subsection{Results}

The first two results are upper bounds obtained from random colourings.
Logarithms are to base $2$.

\begin{theorem}
\label{thm:f}
For every $n\ge1$,
\[
f(n)\ \le\ n^2+n+1
\qquad\text{and}\qquad
f(n)\ <\ 12\,n\,(\log_2(n+2))^2 .
\]
\end{theorem}

\begin{theorem}
\label{thm:F}
Let $n\ge1$ and let $d\ge1$ be an integer with $2^d>nd+2$. Then
\[
F(n)\ \le\ \sum_{j=0}^{d-1}\binom nj .
\]
In particular, for $n\ge32$ one may take
$d=\lceil\log_2n+\log_2\log_2n\rceil+1$, and then
\[
F(n)\ \le\ n^{\,\log_2n+\log_2\log_2n+1}
=\exp\bigl((1+o(1))(\log_2 n)(\ln n)\bigr),
\]
so that $F(n)^{1/n}\to1$: the second conjecture of Erd\H{o}s holds, in
the strong form $F(n)\le\exp(O((\log n)^2))$.
\end{theorem}

Theorem~\ref{thm:F} together with \eqref{eq:trivial} places $F(n)$
between $(n+1)/2$ and $n^{(1+o(1))\log_2n}$. The next result shows that
the random colouring, which is the only tool in
Theorems~\ref{thm:f} and~\ref{thm:F}, cannot close this gap to a
polynomial, and cannot bring $f(n)$ below $n$.

\begin{theorem}
\label{thm:limits}
Let $\chi$ be a uniformly random $2$-colouring of $2^{[n]}$.
\begin{enumerate}[label=\textup{(\roman*)},leftmargin=2em]
\item With probability at least $1/2$ there is a monochromatic chain of
at least $n$ sets; consequently, with probability at least $1/2$, there is
a monochromatic sublattice of size at least $n$.
\item The expected size of the largest monochromatic union-closed family
is at least
\[
\sum_{j=0}^{n}\binom nj2^{1-2^j}\ \ge\ \frac2n\Bigl(\frac n{j_0}\Bigr)^{j_0}
\qquad(n\ge4,\ j_0=\lfloor\log_2\log_2n\rfloor),
\]
and hence at least $n^{\log_2\log_2n-2-o(1)}$.
\end{enumerate}
\end{theorem}

The exponent $\log_2\log_2n$ in (ii) and the exponent $\log_2n$ in
Theorem~\ref{thm:F} are the two sides of the same phenomenon: the
union-closed families that a random colouring cannot avoid are the
complements of the down-sets $\{Y:\chi([n]\setminus Y')=\chi([n])
\text{ for all }Y'\subseteq Y\}$, which have Vapnik--Chervonenkis dimension
about $\log_2\log_2n$, while the Sauer--Shelah bound used in
Theorem~\ref{thm:F} only controls the dimension up to about $\log_2n$.

The last result is computational and concerns small $n$.

\begin{theorem}
\label{thm:small}
The exact values of $f(n)$ and $F(n)$ for small $n$ are as follows.
\[
\begin{array}{c|cccccccccc}
n & 1&2&3&4&5&6&7&8&9&10\\ \hline
f(n) & 1&2&2&3&3&4&4&5&5&6\\
F(n) & 1&2&2&3&4&5&7&&&
\end{array}
\]
In particular $f(n)=\lceil(n+1)/2\rceil$ for $1\le n\le10$, so that the
trivial bound \eqref{eq:trivial} is sharp for $f$ in this range. For each
$n\le10$ there is a colouring attaining $f(n)$ that depends only on the
cardinality of the set; such colourings are listed in
Section~\ref{sec:small}.
\end{theorem}

The computations use a counterexample-guided satisfiability search, described
in Section~\ref{sec:small}, and the program is supplied with the paper.

\subsection{Organisation}
Section~\ref{sec:prelim} collects the structure theory of sublattices of
$2^{[n]}$ (join-irreducibles, Birkhoff's correspondence with down-sets,
Dilworth's theorem) and of union-closed families (independent systems,
Vapnik--Chervonenkis dimension, the Sauer--Shelah lemma).
Section~\ref{sec:f} proves Theorem~\ref{thm:f}, Section~\ref{sec:F}
proves Theorem~\ref{thm:F}, Section~\ref{sec:limits} proves
Theorem~\ref{thm:limits}, Section~\ref{sec:small} describes the
computations behind Theorem~\ref{thm:small}, Section~\ref{sec:questions}
lists the questions that remain, and Section~\ref{sec:lean} describes the
Lean formalisation.

\subsection{Related work}
\label{sec:related}

The problem studied here is easy to confuse with a more famous one that it
resembles only superficially, so the distinction is worth making precise.
\emph{Frankl's union-closed sets conjecture} \cite{Fr79}, posed in 1979,
concerns a single union-closed family $\cF\ne\{\emptyset\}$ of finite sets
and asks whether some element lies in at least half the members of $\cF$;
it is a statement about the internal structure of one family, with no
colouring anywhere in it. It was open in general until Gilmer
\cite{Gi22} proved in 2022, by an entropy argument, that some
element lies in at least a fixed positive proportion (he obtained
$0.01$) of the members of every union-closed family; the constant was
improved to $0.38234$ by Alweiss, Huang and Sellke \cite{AHS24} shortly
afterwards, with $0.5$, the full conjecture, still open at the time of
writing. The problem of the present paper is a Ramsey-type question about
\emph{every} $2$-colouring of the entire Boolean lattice $2^{[n]}$: it
asks not whether a given union-closed family is internally balanced, but
how large a monochromatic union-closed (or intersection-and-union-closed)
family a colouring is forced to contain. The two problems share only
their two words ``union-closed''; a positive answer to one has no direct
bearing on the other, and the entropy method of \cite{Gi22,AHS24}, which
compares a family to its own down-set, does not seem to have an analogue
here, since there is no single family to compare a colour class to before
the colour class itself is fixed.

The Sauer--Shelah lemma (Lemma~\ref{lem:indep}(c)) originates in extremal
set theory \cite{Sa72,Sh72} and is the combinatorial core of the theory of
Vapnik--Chervonenkis dimension \cite{VC71} in statistical learning theory
and of NIP model theory; its use here, bounding a union-closed family by
the independent systems it contains, is close to its original
combinatorial form rather than to its statistical applications, and the
textbook treatments in \cite[Chapter 17]{Bo86} and \cite[Section
3.4]{St12} cover both the lemma and Dilworth's theorem \cite{Di50} used in
Lemma~\ref{lem:posets}. Erd\H{o}s's own problem list \cite{Er78} groups
the question of this paper with several other coloured extremal
set-theory problems from the same period; Bloom's continuously updated
database of Erd\H{o}s's problems \cite{Bl26} records the second conjecture
resolved in Theorem~\ref{thm:F} and the first, left open here, as
Problem~1183, and is the source of the precise attributions used
throughout.

\section{Preliminaries}
\label{sec:prelim}

\subsection{Sublattices of the Boolean lattice}

Throughout this subsection $\cL\subseteq2^{[n]}$ is a sublattice with
$|\cL|\ge2$, that is, a nonempty family closed under union and
intersection. Put
\[
B_0=\bigcap_{A\in\cL}A\in\cL,\qquad T_0=\bigcup_{A\in\cL}A\in\cL .
\]
A member $J\in\cL$ is \emph{join-irreducible} if $J\ne B_0$ and $J$ is
not the union of the members of $\cL$ properly contained in $J$. Let
$\cJ(\cL)$ be the set of join-irreducible members and $k=|\cJ(\cL)|$.
For $x\in T_0\setminus B_0$ put
\[
A_x=\bigcap\{A\in\cL:\ x\in A\}\in\cL ,
\]
the smallest member of $\cL$ containing $x$. A \emph{down-set} of a
finite poset $P$ is a subset $D$ with $y\le x\in D\Rightarrow y\in D$;
$\Down(P)$ denotes the set of down-sets. The \emph{width} $\width(P)$ is the
largest size of an antichain.

\begin{lemma}[Structure of sublattices]
\label{lem:structure}
Let $\cL$ be a sublattice of $2^{[n]}$ with $|\cL|\ge2$.
\begin{enumerate}[label=\textup{(\alph*)},leftmargin=2em]
\item For $x\in T_0\setminus B_0$, $A_x$ is join-irreducible, and every
join-irreducible member is of the form $A_x$.
\item Every $A\in\cL$ equals $B_0\cup\bigcup\{J\in\cJ(\cL):J\subseteq A\}$.
\item If $J\in\cJ(\cL)$ and $J\subseteq A\cup A'$ with $A,A'\in\cL$, then
$J\subseteq A$ or $J\subseteq A'$. The same holds for unions of any
number of members.
\item The map $A\mapsto D(A)=\{J\in\cJ(\cL):J\subseteq A\}$ is a
bijection from $\cL$ onto $\Down(\cJ(\cL),\subseteq)$; in particular
$|\cL|=|\Down(\cJ(\cL),\subseteq)|$.
\item $k=|\cJ(\cL)|\le n$, and $\width(\cJ(\cL),\subseteq)\le\log_2|\cL|$.
\item Define $\psi:[n]\to\cJ(\cL)\cup\{\bot,\top\}$ by $\psi(x)=\bot$ if
$x\in B_0$, $\psi(x)=\top$ if $x\notin T_0$, and $\psi(x)=A_x$ otherwise.
Then $\cL=\{A_D:D\in\Down(\cJ(\cL),\subseteq)\}$ with
$A_D=\{x:\psi(x)=\bot\text{ or }\psi(x)\in D\}$. Consequently $\cL$ is
determined by the poset $(\cJ(\cL),\subseteq)$, given abstractly on a set of
$k$ labels, together with the map $\psi$.
\end{enumerate}
\end{lemma}

\begin{proof}
(a) $A_x\in\cL$ since $\cL$ is closed under intersection and $T_0$ is a
member containing $x$; $A_x\ne B_0$ since $x\notin B_0$; and if $A_x$ were
the union of members properly contained in it, one of them would contain
$x$, contradicting the minimality of $A_x$. Conversely let $J\in\cJ(\cL)$
and let $U$ be the union of the members of $\cL$ properly contained in
$J$; since $B_0\subseteq J$ and $B_0\ne J$, $U\in\cL$ and $B_0\subseteq
U\subsetneq J$. Pick $x\in J\setminus U$. Then $A_x\subseteq J$, and
$A_x\subsetneq J$ would give $A_x\subseteq U$, which is false; so
$A_x=J$.

(b) Let $A\in\cL$. Then $B_0\subseteq A$, and for $x\in A\setminus B_0$
we have $A_x\subseteq A$ with $A_x\in\cJ(\cL)$ by (a), so $A$ is contained
in the right side; the reverse inclusion is clear.

(c) Since $\cL$ is closed under intersection, $J=J\cap(A\cup A')
=(J\cap A)\cup(J\cap A')$ with $J\cap A,J\cap A'\in\cL$ contained in $J$.
If both were proper subsets of $J$, then $J$ would be the union of two
members properly contained in it, contradicting join-irreducibility.
Hence $J\cap A=J$ or $J\cap A'=J$. For a union of several members apply
this repeatedly.

(d) By (b) the map is injective. It takes values in down-sets: if
$J'\subseteq J\subseteq A$ then $J'\subseteq A$. For surjectivity let $D$
be a down-set and put $A_D=B_0\cup\bigcup_{J\in D}J\in\cL$. If
$J\in\cJ(\cL)$ satisfies $J\subseteq A_D$, then by (c) applied to the
union $B_0\cup\bigcup_{J'\in D}J'$ we have $J\subseteq B_0$, impossible as
$B_0\subsetneq J$, or $J\subseteq J'$ for some $J'\in D$, whence $J\in D$
as $D$ is a down-set. Thus $D(A_D)=D$.

(e) Order $\cJ(\cL)$ by a linear extension $J_1,\dots,J_k$ of $\subseteq$;
the sets $D_i=\{J_1,\dots,J_i\}$, $0\le i\le k$, are down-sets, so by (d)
$A_{D_0}\subsetneq A_{D_1}\subsetneq\dots\subsetneq A_{D_k}$ is a chain
of $k+1$ distinct subsets of $[n]$, whence $k+1\le n+1$. If $S$ is an
antichain of $\cJ(\cL)$, the $2^{|S|}$ sets $\{J\in\cJ(\cL):J\subseteq
J'\text{ for some }J'\in S'\}$, $S'\subseteq S$, are distinct down-sets,
so $2^{\width}\le|\cL|$.

(f) By (d) every member is $A_D=B_0\cup\bigcup_{J\in D}J$ for a unique
down-set $D$. An element $x\in[n]$ lies in $A_D$ if and only if $x\in B_0$,
or $x\in T_0\setminus B_0$ and $x\in J$ for some $J\in D$; and for
$x\in T_0\setminus B_0$, $x\in J$ holds if and only if $A_x\subseteq J$,
by minimality of $A_x$. Since $D$ is a down-set, $A_x\subseteq J\in D$ for
some $J$ if and only if $A_x\in D$. This is the displayed description of
$A_D$, which uses only $\psi$ and the order structure of $\cJ(\cL)$.
\end{proof}

Part (d) is Birkhoff's representation of finite distributive lattices
\cite[Theorem 3.4.1]{St12}, specialised to sublattices of $2^{[n]}$; we have
included the short proof since (c) and (f) are used in exactly this form.

\begin{lemma}[Counting posets of bounded width]
\label{lem:posets}
The number of partial orders on the labelled set $[k]$ of width at most
$w$ is at most $(k+1)^{k(w+2)}$.
\end{lemma}

\begin{proof}
Let $P$ be such a poset. By Dilworth's theorem \cite{Di50}, $[k]$ is the
disjoint union of $w$ chains $C_1,\dots,C_w$ of $P$, some possibly empty.
Record the following data: for each $x\in[k]$ the index $c(x)\in[w]$ of
its chain and its position $p(x)\in[k]$ within that chain (counted from the
bottom); and for each $x\in[k]$ and each $i\in[w]$ the number
$r_i(x)\in\{0,\dots,k\}$ of elements $y\in C_i$ with $y\le x$. Since
$y\le x$ and $y'\le y$ with $y'\in C_i$ imply $y'\le x$, the elements of
$C_i$ below $x$ form an initial segment of $C_i$, and so $y\le x$ holds if
and only if $p(y)\le r_{c(y)}(x)$. Thus the data determine $P$, and the
number of possible data is at most $w^k\cdot k^k\cdot(k+1)^{kw}\le
(k+1)^{k(w+2)}$, using $w\le k$.
\end{proof}

\begin{lemma}[Counting sublattices]
\label{lem:count}
For $m\ge2$, the number $N(m)$ of sublattices $\cL\subseteq2^{[n]}$ with
$|\cL|=m$ satisfies
\[
N(m)\ \le\ \sum_{k=1}^{n}(k+1)^{k(\lfloor\log_2m\rfloor+2)}(k+2)^n
\ \le\ (n+2)^{\,n(\log_2m+4)} .
\]
Moreover the total number of sublattices of $2^{[n]}$ with at least two
members is at most $2^{n(n+1)}$.
\end{lemma}

\begin{proof}
By Lemma~\ref{lem:structure}(f), $\cL$ is determined by a poset on $k$
labels together with a map $\psi:[n]\to[k]\cup\{\bot,\top\}$, where
$k\le n$ and the poset has width at most $\log_2m$ by
Lemma~\ref{lem:structure}(e). Lemma~\ref{lem:posets} bounds the number of
posets and $(k+2)^n$ bounds the number of maps, which gives the first
inequality. For the second, each summand is at most
$(n+1)^{n(\log_2m+2)}(n+2)^n\le(n+2)^{n(\log_2m+2)+n}$, and there are
$n\le(n+2)^n$ summands. For the last assertion, $\cL$ is determined by
$B_0$ together with $\cJ(\cL)$, by Lemma~\ref{lem:structure}(b) and the
fact that $\cL$ consists of all unions $B_0\cup\bigcup_{J\in D}J$; this is
a list of at most $n+1$ subsets of $[n]$, and there are at most
$(2^n)^{n+1}$ such lists.
\end{proof}

\subsection{Union-closed families}

\begin{definition}
Sets $A_1,\dots,A_d\subseteq[n]$ form an \emph{independent system} if
each $A_i$ has a \emph{private element}: an $x_i\in A_i$ with
$x_i\notin A_j$ for all $j\ne i$. The \emph{independence number}
$\ind(\cF)$ of a family $\cF$ is the largest $d$ for which $\cF$ contains an
independent system of $d$ sets. A set $X\subseteq[n]$ is \emph{shattered}
by $\cF$ if $\{A\cap X:A\in\cF\}=2^X$, and the
\emph{Vapnik--Chervonenkis dimension} $\vc(\cF)$ is the largest size of a
shattered set.
\end{definition}

\begin{lemma}
\label{lem:indep}
\begin{enumerate}[label=\textup{(\alph*)},leftmargin=2em]
\item If $A_1,\dots,A_d$ is an independent system, then the $2^d-1$
unions $\bigcup_{i\in I}A_i$, $\emptyset\ne I\subseteq[d]$, are pairwise
distinct.
\item For every family $\cF$, $\vc(\cF)\le\ind(\cF)$.
\item \textup{(Sauer \cite{Sa72}, Shelah \cite{Sh72}.)} If
$\vc(\cF)\le d$ then $|\cF|\le\sum_{j=0}^{d}\binom nj$.
\end{enumerate}
\end{lemma}

\begin{proof}
(a) The private element $x_i$ lies in $\bigcup_{i\in I}A_i$ if and only if
$i\in I$, so the union determines $I$.
(b) Let $X=\{x_1,\dots,x_d\}$ be shattered, and for each $i$ choose
$A_i\in\cF$ with $A_i\cap X=\{x_i\}$. Then $x_i$ is a private element of
$A_i$ in the system $A_1,\dots,A_d$.
(c) This is the Sauer--Shelah lemma.
\end{proof}

\section{Upper bounds for $f(n)$}
\label{sec:f}

Let $\chi$ be a uniformly random $2$-colouring of $2^{[n]}$: the colours
$\chi(A)$, $A\subseteq[n]$, are independent and uniform on $\{0,1\}$.
A fixed family of $m$ sets is monochromatic with probability $2^{1-m}$.

\begin{proof}[Proof of Theorem~\ref{thm:f}, first bound]
Let $M=n^2+n+2$. By Lemma~\ref{lem:count} there are at most $2^{n(n+1)}$
sublattices with at least two members, so the expected number of
monochromatic sublattices with at least $M$ members is at most
$2^{n(n+1)}\cdot2^{1-M}=2^{-1}<1$. Hence some colouring has no
monochromatic sublattice with $M$ or more members, and $f(n)\le M-1$.
\end{proof}

\begin{proof}[Proof of Theorem~\ref{thm:f}, second bound]
For $n\le2$ the bound is weaker than the trivial $f(n)\le2^n$, so let
$n\ge3$ and put $\ell=\log_2(n+2)\ge\log_25>2.32$, $A=n\ell$ and
$M=\lceil12n\ell^2\rceil$. By Lemma~\ref{lem:count}, for every $m\ge2$,
\[
N(m)\,2^{1-m}\le2^{A(\log_2m+4)+1-m}.
\]
Let $\varphi(m)=m/2-A(\log_2m+4)-1$. We claim $\varphi(m)\ge0$ for all
real $m\ge M$. The derivative $\varphi'(m)=1/2-A/(m\ln2)$ is positive for
$m>2A/\ln2$, and $M\ge12n\ell^2>2.9\,n\ell>2A/\ln2$, so it suffices to
check $\varphi(M)\ge0$. Using $M\le13n\ell^2$ and
$\log_2\ell\le\ell-1$ (valid for $\ell\ge1$),
\[
\log_2M+4\le\log_213+\log_2n+2\log_2\ell+4
\le3.71+\ell+2(\ell-1)+4=3\ell+5.71,
\]
so $A(\log_2M+4)+1\le3n\ell^2+5.71n\ell+1\le6n\ell^2\le M/2$, the middle
inequality being $3\ell\ge5.71+1/(n\ell)$, which holds as $\ell>2.32$.
Hence, for $m\ge M$, $N(m)2^{1-m}\le2^{-m/2}$, and the expected number of
monochromatic sublattices with at least $M$ members is at most
$\sum_{m\ge M}2^{-m/2}<2^{-M/2}/(1-2^{-1/2})<1$. Some colouring therefore
has no monochromatic sublattice with $M$ or more members, and
$f(n)\le M-1<12n\ell^2$.
\end{proof}

\begin{remark}
The first bound is the better one for $n\le1500$ or so, the second
thereafter. Both use only the number of sublattices and the probability
$2^{1-m}$; the second is limited by the factor $(n+2)^n$ in
Lemma~\ref{lem:count}, which is really present (there are $n!$ maximal
chains), so this method cannot give $f(n)=O(n)$. Theorem~\ref{thm:limits}(i)
shows that the random colouring itself has monochromatic chains of $n$
sets, so no analysis of the random colouring can bring the bound below
$n$.
\end{remark}

\section{Upper bound for $F(n)$}
\label{sec:F}

\begin{proof}[Proof of Theorem~\ref{thm:F}]
Let $\chi$ be a uniformly random $2$-colouring and let $d\ge1$ satisfy
$2^d>nd+2$. Call an ordered $d$-tuple $(A_1,\dots,A_d)$ of subsets of
$[n]$ \emph{bad} if it is an independent system and all $2^d-1$ unions
$\bigcup_{i\in I}A_i$, $\emptyset\ne I\subseteq[d]$, have the same colour.
By Lemma~\ref{lem:indep}(a) these unions are distinct sets, so a fixed
independent tuple is bad with probability $2^{1-(2^d-1)}=2^{2-2^d}$.
There are at most $2^{nd}$ tuples, so the expected number of bad tuples
is at most $2^{nd+2-2^d}<1$. Fix a colouring with no bad tuple.

Let $\cF$ be a monochromatic union-closed family for this colouring. If
$\cF$ contained an independent system $A_1,\dots,A_d$, then all unions
$\bigcup_{i\in I}A_i$ would belong to $\cF$ by union-closure and would
therefore be monochromatic, making the tuple bad. Hence $\ind(\cF)\le d-1$,
so $\vc(\cF)\le d-1$ by Lemma~\ref{lem:indep}(b), and
$|\cF|\le\sum_{j\le d-1}\binom nj$ by Lemma~\ref{lem:indep}(c). This
proves the first assertion.

For the second, let $n\ge32$ and $d=\lceil\log_2n+\log_2\log_2n\rceil+1$.
Then $2^d\ge2n\log_2n$, while $nd+2\le n(\log_2n+\log_2\log_2n+2)+2$; the
inequality $2n\log_2n>n\log_2n+n\log_2\log_2n+2n+2$ is equivalent to
$\log_2n-\log_2\log_2n>2+2/n$, which holds for $n\ge32$ (at $n=32$ the
left side is $5-\log_25>2.67$ and the left side is increasing). Finally
$\sum_{j\le d-1}\binom nj\le n^{d-1}\le n^{\log_2n+\log_2\log_2n+1}$, and
$n^{\log_2n+\log_2\log_2n+1}=\exp((\log_2n+\log_2\log_2n+1)\ln n)$, whose
$n$-th root tends to $1$.
\end{proof}

\begin{remark}
For $n<32$ the theorem applies with other values of $d$: $d=6$ for
$n\le10$, $d=7$ for $n\le18$, $d=8$ for $n\le31$. The resulting bounds are
far from the true values in Theorem~\ref{thm:small}, as they must be, since
the method does not see the chain structure at all.
\end{remark}

\begin{remark}[An explicit correction]
\label{rem:explicit}
The inequality $\sum_{j\le d-1}\binom nj\le n^{d-1}$ used at the end of
the proof of Theorem~\ref{thm:F} is not quite literally true: since
$\binom nj\le n^j$ for every $j\ge0$ and $n\ge1$, one has only
\[
\sum_{j=0}^{d-1}\binom nj\ \le\ \sum_{j=0}^{d-1}n^j\ \le\ d\,n^{d-1}
\qquad(n\ge1),
\]
the last step bounding each of the $d$ terms by the largest, $n^{d-1}$.
This is the correct form of the estimate, and it changes nothing: with
$d=\lceil\log_2n+\log_2\log_2n\rceil+1=O(\log n)$, the extra factor $d$
contributes $\log_2d=O(\log\log n)$ to $\log_2$ of the bound, which is
already absorbed by the $o(1)$ in the exponent $\log_2n+\log_2\log_2n+1$.
Precisely,
\begin{equation}
\label{eq:explicit}
F(n)\ \le\ d\cdot n^{d-1}\ \le\ (\log_2n+\log_2\log_2n+2)\cdot
n^{\,\log_2n+\log_2\log_2n+1}\qquad(n\ge32),
\end{equation}
which is the fully justified version of the last display in the proof of
Theorem~\ref{thm:F} and is the bound tabulated in
Table~\ref{tab:explicit} below.
\end{remark}

\section{Explicit and numerical bounds}
\label{sec:explicit}

Theorem~\ref{thm:F} is stated for $n\ge32$ so that the two-line
computation in its proof goes through without case analysis; the
underlying first-moment argument needs no such threshold, and the
following restatement removes it.

\begin{proposition}[Threshold-free form]
\label{prop:allN}
For every $n\ge1$, let $d(n)$ be the least positive integer $d$ with
$2^d>nd+2$. Then
\[
F(n)\ \le\ \sum_{j=0}^{d(n)-1}\binom nj\ \le\ d(n)\,n^{d(n)-1}.
\]
The function $d(n)$ is nondecreasing, takes the values
$3,4,4,5,5,6,6,6,6,6$ for $n=1,\dots,10$, and satisfies
$d(n)=\log_2n+\log_2\log_2n+O(1)$ as $n\to\infty$; in particular
$d(n)/\log_2n\to1$, so $F(n)\le n^{(1+o(1))\log_2n}$ and $F(n)^{1/n}\to1$
with no restriction on $n$.
\end{proposition}

\begin{proof}
The first inequality is the first assertion of Theorem~\ref{thm:F},
proved for every $n\ge1$ and every $d$ with $2^d>nd+2$ (the proof uses
only this hypothesis, not $n\ge32$); taking $d=d(n)$, the least such $d$,
gives the strongest instance. The second inequality is
Remark~\ref{rem:explicit}. Monotonicity of $d(n)$ holds because $2^d>nd+2$
implies $2^d>n'd+2$ for $n'\le n$. The values for $n\le9$ and the
asymptotics for $d(n)$ are immediate from the defining inequality: writing
$d=\log_2n+\log_2\log_2n+e(n)$, the condition $2^d>nd+2$ becomes, after
dividing by $n\log_2n$, an inequality in $e(n)$ that is satisfied once
$e(n)\ge2+o(1)$ and fails once $e(n)\le-o(1)$, giving $e(n)=O(1)$; this is
the computation carried out for $n\ge32$ in the proof of
Theorem~\ref{thm:F} and it is valid, with different explicit constants,
for every $n$.
\end{proof}

Table~\ref{tab:explicit} evaluates the bound of Proposition~\ref{prop:allN}
directly: for each $n$ it records $d(n)$, the exact value of
$\log_2\sum_{j<d(n)}\binom nj$ for $n\le2000$ (computed by exact integer
arithmetic) or the bound $\log_2(d(n))+(d(n)-1)\log_2n$ of
\eqref{eq:explicit} for larger $n$, and the ratio of this quantity to $n$,
which is $\log_2(F(n)^{1/n})$ up to the direction of the inequality and
which Theorem~\ref{thm:F} and Proposition~\ref{prop:allN} assert tends to
$0$.

\begin{table}[h]
\centering
\begin{tabular}{r r r r}
\toprule
$n$ & $d(n)$ & $\log_2(\text{bound})$ & $\log_2(\text{bound})/n$\\
\midrule
$10$ & $6$ & $9.317$ & $0.9317$\\
$32$ & $9$ & $23.842$ & $0.7451$\\
$50$ & $9$ & $29.287$ & $0.5857$\\
$100$ & $10$ & $40.937$ & $0.4094$\\
$200$ & $12$ & $58.514$ & $0.2926$\\
$500$ & $13$ & $78.598$ & $0.1572$\\
$1000$ & $14$ & $96.925$ & $0.0969$\\
$10^4$ & $18$ & $225.891$ & $0.02259$\\
$10^5$ & $22$ & $348.802$ & $0.00349$\\
$10^6$ & $25$ & $478.358$ & $0.00048$\\
\bottomrule
\end{tabular}
\caption{The bound of Proposition~\ref{prop:allN}; entries for
$n\le2000$ use the exact value of $\log_2\sum_{j<d(n)}\binom nj$, entries
for $n>2000$ use the right side of \eqref{eq:explicit}. The last column,
which is at least $\log_2(F(n))/n$ by Theorem~\ref{thm:F} and at most the
value shown, illustrates the convergence $F(n)^{1/n}\to1$ at concrete
scales; it says nothing about the true order of $F(n)$, which by
Theorem~\ref{thm:limits}(ii) is not polynomial along the random colouring
and is addressed only for cardinality colourings, in
Section~\ref{sec:card}.}
\label{tab:explicit}
\end{table}

\section{Limits of the random colouring}
\label{sec:limits}

\begin{proof}[Proof of Theorem~\ref{thm:limits}]
(i) Build a chain greedily. Put $S_0=\emptyset$ and $c=\chi(\emptyset)$.
Given $S_i$ with $|S_i|=i<n$ and $\chi(S_i)=c$, the $n-i$ sets
$S_i\cup\{x\}$, $x\notin S_i$, have independent uniform colours, none of
which has been examined before (each is examined only at the step at
which it is a candidate, and $S_i$ determines the step). The step fails
if all of them have colour $1-c$, which has probability $2^{-(n-i)}$;
otherwise choose one of colour $c$ as $S_{i+1}$. The probability that no
step fails for $i=0,\dots,n-2$ is at least
$1-\sum_{i=0}^{n-2}2^{-(n-i)}=1-\sum_{j=2}^{n}2^{-j}>1/2$, and in that
event $S_0\subset S_1\subset\dots\subset S_{n-1}$ is a monochromatic
chain of $n$ sets, which is a sublattice.

(ii) Let $T=[n]$ and $c=\chi(T)$, and let
\[
\cY=\{Y\subseteq T:\ \chi(T\setminus Y')=c\text{ for all }Y'\subseteq Y\}.
\]
Then $\cY$ is a down-set of $2^{[n]}$ containing $\emptyset$, and the
family $\cF=\{T\setminus Y:Y\in\cY\}$ is monochromatic of colour $c$ and
union-closed: $(T\setminus Y)\cup(T\setminus Y')=T\setminus(Y\cap Y')$ and
$Y\cap Y'\in\cY$. A set $Y$ with $|Y|=j$ lies in $\cY$ if and only if the
$2^j-1$ sets $T\setminus Y'$, $\emptyset\ne Y'\subseteq Y$, all have the
colour of $T$, an event of probability $2^{-(2^j-1)}$. Hence
$\Ex|\cF|=\Ex|\cY|=\sum_{j=0}^{n}\binom nj2^{1-2^j}$, and the largest
monochromatic union-closed family has expected size at least this. For
the numerical bound take $j_0=\lfloor\log_2\log_2n\rfloor$ with $n\ge4$;
then $2^{j_0}\le\log_2n$, so $2^{1-2^{j_0}}\ge2/n$, and
$\binom n{j_0}\ge(n/j_0)^{j_0}$; thus $\Ex|\cF|\ge2(n/j_0)^{j_0}/n$. As
$(n/j_0)^{j_0}=n^{j_0-o(1)}$ and $j_0\ge\log_2\log_2n-1$, this is
$n^{\log_2\log_2n-2-o(1)}$.
\end{proof}

\begin{remark}
Part (ii) shows that no bound of the form $F(n)\le n^{C}$ with a constant
$C$ can be proved by exhibiting a random colouring; the same family also
shows that for the random colouring the Vapnik--Chervonenkis dimension of
the largest monochromatic union-closed family is at least about
$\log_2\log_2n$, since $\cY$ contains all sets of size at most $j_0$ with
positive probability bounded away from zero when $2^{j_0}\le\log_2n$.
\end{remark}


\section{Colourings depending only on cardinality}
\label{sec:card}

A colouring $\chi$ is a \emph{cardinality colouring} if $\chi(S)=g(|S|)$
for some $g:\{0,1,\dots,n\}\to\{0,1\}$. These are the colourings for which
Erd\H{o}s recorded Howorka's assertion, and they are the colourings that
attain $f(n)$ for $n\le10$ (Section~\ref{sec:small}). For a cardinality
colouring, monochromatic closed families can be built from monochromatic
arithmetic progressions of $g$.

\begin{proposition}
\label{prop:AP}
Let $\chi(S)=g(|S|)$ and suppose that the arithmetic progression
$a,a+s,\dots,a+(k-1)s$, with $s\ge1$, $k\ge1$ and $a+(k-1)s\le n$, is
monochromatic for $g$. Put $q=\lfloor a/s\rfloor$ and $b=q+k-1$. Then
there is a monochromatic union-closed family with
$\sum_{i=0}^{k-1}\binom bi$ members, and a monochromatic sublattice with
$2^{k-1}$ members. More generally, if $t_1<t_2<\dots<t_L$ and $s,w\ge1$
are such that $t_L+ws\le n$ and all the numbers $t_j+is$
($1\le j\le L$, $0\le i\le w$) have the same colour under $g$, then there
is a monochromatic sublattice with $L\,2^{w}$ members.
\end{proposition}

\begin{proof}
Put $a_0=a-qs\in[0,s)$ and choose pairwise disjoint sets
$B,X_1,\dots,X_b\subseteq[n]$ with $|B|=a_0$ and $|X_i|=s$; this is
possible since $a_0+bs=a+(k-1)s\le n$. Let
\[
\cF=\Bigl\{B\cup\bigcup_{i\in I}X_i:\ I\subseteq[b],\ |I|\ge q\Bigr\}.
\]
A member with index set $I$ has cardinality $a_0+|I|s=a+(|I|-q)s$ with
$0\le|I|-q\le k-1$, so all members have the colour of the progression.
The union of the members with index sets $I,I'$ has index set
$I\cup I'$, of size at least $q$, so $\cF$ is union-closed, and
$|\cF|=\sum_{i=q}^{b}\binom bi=\sum_{i=0}^{k-1}\binom bi$. For the
sublattice take $|B'|=a$ and $k-1$ disjoint blocks $X_1,\dots,X_{k-1}$ of
size $s$ disjoint from $B'$, and the family of all $B'\cup\bigcup_{i\in I}
X_i$, $I\subseteq[k-1]$; it is closed under union and intersection, has
$2^{k-1}$ members, and their cardinalities $a+|I|s$ lie in the
progression.

For the last assertion choose disjoint blocks $X_1,\dots,X_w$ of size $s$
and a chain $C_1\subset\dots\subset C_L$ of subsets of
$[n]\setminus(X_1\cup\dots\cup X_w)$ with $|C_j|=t_j$, possible since
$t_L\le n-ws$. The family $\{C_j\cup\bigcup_{i\in I}X_i:\ 1\le j\le L,\
I\subseteq[w]\}$ has $L2^w$ members of cardinalities $t_j+|I|s$, all of
one colour, and it is closed under union and intersection since
$(C_j\cup X_I)\cup(C_{j'}\cup X_{I'})=C_{\max(j,j')}\cup X_{I\cup I'}$ and
$(C_j\cup X_I)\cap(C_{j'}\cup X_{I'})=C_{\min(j,j')}\cup X_{I\cap I'}$.
\end{proof}

Let $W(k)$ denote the van der Waerden number: the least $N$ such that
every $2$-colouring of $\{1,\dots,N\}$ contains a monochromatic arithmetic
progression of $k$ terms; it is finite by van der Waerden's theorem
\cite{vdW27}. The following statement is the one attributed to Howorka in
\cite{Er78}; we have not been able to locate a published proof, and give
one.

\begin{theorem}
\label{thm:howorka}
Let $k\ge2$ and $n\ge2W(k)$. For every cardinality colouring of
$2^{[n]}$ there is a monochromatic union-closed family with at least
\[
\binom{\lfloor n/(2W(k))\rfloor+k-1}{k-1}\ \ge\
\Bigl(\frac{n}{2(k-1)W(k)}\Bigr)^{k-1}
\]
members, and a monochromatic sublattice with $2^{k-1}$ members.
Consequently, for every constant $c$, every cardinality colouring of
$2^{[n]}$ with $n$ large in terms of $c$ has a monochromatic union-closed
family of more than $n^{c}$ members; and there is a function
$\omega(n)\to\infty$ such that every cardinality colouring has a
monochromatic union-closed family of more than $n^{\omega(n)}$ members.
\end{theorem}

\begin{proof}
Let $\chi(S)=g(|S|)$ and $W=W(k)$. The $W$ consecutive integers
$m_0,m_0+1,\dots,m_0+W-1$ with $m_0=\lceil n/2\rceil$ lie in
$\{0,\dots,n\}$ because $n\ge2W$, and by the definition of $W(k)$ the
colouring $g$ restricted to them contains a monochromatic progression
$a,a+s,\dots,a+(k-1)s$ with $a\ge n/2$ and $(k-1)s\le W-1$, so that
$s<W$ and $\lfloor a/s\rfloor\ge\lfloor n/(2W)\rfloor$. By
Proposition~\ref{prop:AP} there is a monochromatic union-closed family
with $\sum_{i<k}\binom bi\ge\binom b{k-1}$ members, where
$b=\lfloor a/s\rfloor+k-1\ge\lfloor n/(2W)\rfloor+k-1$, and
$\binom b{k-1}\ge(b/(k-1))^{k-1}\ge(n/(2(k-1)W))^{k-1}$, using
$b\ge\lfloor n/(2W)\rfloor+k-1\ge n/(2W)$ (as $k\ge2$).
The sublattice statement is the second part of Proposition~\ref{prop:AP}.
For fixed $k$ the bound is $c_kn^{k-1}$ with $c_k>0$, which exceeds
$n^{c}$ for $n$ large once $k-1>c$. Finally let $\omega(n)$ be the largest
$k$ with $2(k-1)W(k)\le n^{1/2}$; then $\omega(n)\to\infty$, and the bound
is at least $n^{(\omega(n)-1)/2}$, which exceeds $n^{\omega(n)/3}$ for
$\omega(n)\ge3$; replacing $\omega$ by $\omega/3$ gives the statement.
\end{proof}

\begin{remark}
The van der Waerden numbers grow so fast that Theorem~\ref{thm:howorka}
gives nothing for small $n$; its content is the polynomial growth for
every fixed $k$. Proposition~\ref{prop:AP} also explains the shape of the
extremal cardinality colourings of Section~\ref{sec:small}: a cardinality
colouring with no monochromatic sublattice of more than
$\lceil(n+1)/2\rceil$ members must have no monochromatic progression of
$k$ terms with $2^{k-1}>\lceil(n+1)/2\rceil$, and, for every $s$ and $w$,
at most $\lceil(n+1)/2\rceil2^{-w}$ starting points $t$ of one colour with
$t,t+s,\dots,t+ws$ all of that colour. For $n\le10$ both conditions leave
room, as the colourings $g_n$ listed in Section~\ref{sec:small} show.
\end{remark}

\section{Exact values for small $n$}
\label{sec:small}

\subsection{The algorithm}
Fix $n$ and a target $M$. We wish to decide whether there is a
$2$-colouring of $2^{[n]}$ with no monochromatic closed family of $M$ or
more members, closure meaning union and intersection for $f$ and union
alone for $F$. This is a satisfiability problem in the $2^n$ Boolean
variables $x_S$, $S\subseteq[n]$, with one clause pair for every closed
family $\cF$ of size $M$: not all $x_S$, $S\in\cF$, true, and not all
false. Rather than enumerating all closed families in advance, the program
proceeds by counterexample-guided refinement: a SAT solver (CaDiCaL,
through the PySAT interface \cite{IMM18}) proposes a colouring subject to
the clauses collected so far; a branch-and-bound search then looks, in
each colour class, for a closed subfamily with at least $M$ members,
building families by adding sets one at a time and closing under the
operations; if one is found, its clause pair is added and the solver is
run again; if none is found, the colouring is a witness that the value is
less than $M$; if the solver reports unsatisfiability, every colouring
has a monochromatic closed family of at least $M$ members and the value
is at least $M$. The symmetry $\chi\mapsto1-\chi$ is broken by fixing the
colour of $\emptyset$. Increasing $M$ from $1$ gives the exact value.

\subsection{Results and checks}
The values in Theorem~\ref{thm:small} were obtained in this way. The
computation for $f(9)$ examined $14\,190$ clauses and took
twenty minutes on one processor core; the computation for $F(7)$ examined $256\,428$ clauses. Three
independent checks were made. For $n\le4$ all $2^{2^n}$ colourings were
enumerated directly and the maximal monochromatic closed family was
found by inspecting all closed families, without any SAT solver; the
values agree. For every entry of the table the final colouring produced
by the solver (the witness that the value is not larger) was verified by
a second program that encodes ``this colour class contains a closed
family of at least $M$ members'' as a satisfiability instance with a
cardinality constraint, independent of the branch-and-bound search, and
the instance was unsatisfiable for both colour classes in every case;
the witnesses are supplied with the code, and this second, independent
check was re-run for the present version of the paper, confirming for
every $g_n$, $3\le n\le10$, that neither colour class contains a
monochromatic sublattice of $\lceil(n+1)/2\rceil+1$ or more members. In
particular the values of
$f(n)$ for $n\le10$ do not depend on any unsatisfiability report: the
lower bound is \eqref{eq:trivial}, and the upper bound is certified by the
cardinality colourings $g_n$ listed below, each checked by the second
program. The lower bounds for $F(5)$, $F(6)$ and $F(7)$ rest on the
solver's unsatisfiability reports.

The extremal colourings found by the solver for $n\le7$ happened to be
cardinality colourings; for $n=8$ and $n=9$ a cardinality colouring
attaining $f(n)$ was found by adding to the solver the clauses
$x_S\leftrightarrow x_{S'}$ for $|S|=|S'|$, and for $n=10$ (where the
unrestricted search was not run to completion) by enumerating the
cardinality colourings that pass the necessary conditions of
Proposition~\ref{prop:AP} and checking each with the second program;
since $f(10)\ge6$ by \eqref{eq:trivial}, this determines $f(10)=6$. Writing
$g_n=(g_n(0),g_n(1),\dots,g_n(n))$ for the colour of a set of each
cardinality, the following colourings have no monochromatic sublattice
with more than $\lceil(n+1)/2\rceil$ members:
{\small
\begin{align*}
g_3&=(1,1,0,0), & g_4&=(1,1,0,1,0),\\
g_5&=(1,1,0,1,0,0), & g_6&=(1,0,0,1,0,1,1),\\
g_7&=(1,0,0,1,1,0,0,1), & g_8&=(1,0,1,0,0,1,1,0,1),\\
g_9&=(1,0,1,0,0,1,1,0,1,0), & g_{10}&=(1,0,0,1,1,0,0,1,0,1,1).
\end{align*}}
For a colouring depending only on cardinality, a monochromatic sublattice
corresponds, by Lemma~\ref{lem:structure}, to a poset on blocks of
elements together with block sizes and a base size, all of whose down-set
weights lie in one colour class; the reader may check by hand, for
instance, that for $g_9$ the colour class $\{0,2,5,6,8\}$ contains the
arithmetic progression $2,5,8$, which yields the sublattice
$\{\emptyset,B,B\cup X,B\cup Y,B\cup X\cup Y\}$ of five members with
$|B|=2$, $|X|=|Y|=3$, and that this family cannot be extended within the
colour class.

\subsection{Discussion}
The values show that the trivial bound \eqref{eq:trivial} for $f$ is
exact up to $n=10$, and that $F$ exceeds $f$ from $n=5$ on. We do not
conjecture that $f(n)=\lceil(n+1)/2\rceil$ holds for all $n$. By
Proposition~\ref{prop:AP}, a cardinality colouring attaining this value
for all $n$ would have to avoid monochromatic arithmetic progressions of
length $\log_2(n+1)+1$ in every colour class, and to contain, for every
difference $s$ and length $w+1$, at most $\lceil(n+1)/2\rceil2^{-w}$ monochromatic
progressions of that difference and length in each colour; the second
condition is met with room to spare for $n\le10$ and asks, for large $n$,
for a colouring whose progression counts never exceed their expected
values, which we do not expect to exist. Whether $f(n)=O(n)$ is a
different question, since general colourings are not constrained by
Proposition~\ref{prop:AP}.

\section{Summary of what is proved and what is open}
\label{sec:summary}

The results above touch four different quantities, and it is easy to
lose track of which are settled and which are not; this section collects
them in one place.

\begin{table}[h]
\centering
\begin{tabular}{p{0.3\linewidth} p{0.38\linewidth} p{0.22\linewidth}}
\toprule
Question & Status & Where\\
\midrule
Erd\H{o}s's second conjecture, $F(n)<(1+\eps)^n$ for every $\eps>0$ &
PROVED, in the strong form $F(n)\le n^{(1+o(1))\log_2n}$, for every
$n\ge1$ & Theorem~\ref{thm:F}, Prop.~\ref{prop:allN}\\
\addlinespace
Erd\H{o}s's first conjecture, $F(n)>n^c$ for every constant $c$ &
OPEN in general; PROVED for cardinality colourings & Question~3,
Theorem~\ref{thm:howorka}\\
\addlinespace
Order of magnitude of $f(n)$ & OPEN; $f(n)=\lceil(n+1)/2\rceil$ for
$n\le10$, upper bound $f(n)<12n(\log_2n)^2$ in general & Theorems
\ref{thm:f}, \ref{thm:small}\\
\addlinespace
Is $f(n)=O(n)$? & OPEN; the random colouring alone cannot show it,
since it already has monochromatic chains of length $n$ &
Theorem~\ref{thm:limits}(i)\\
\addlinespace
Is $F(n)$ polynomially bounded? & OPEN; known bounds place $\log_2F(n)$
between about $(\log_2n)(\log_2\log_2n)$ and $2\log_2n\cdot\log_2n$ &
Theorems \ref{thm:F}, \ref{thm:limits}(ii)\\
\addlinespace
Exact values of $f(n)$, $F(n)$ for small $n$ & COMPUTED and doubly
verified for $n\le10$ (partially $F(8)$--$F(10)$ open) &
Theorem~\ref{thm:small}, Appendix~\ref{app:sat}\\
\bottomrule
\end{tabular}
\caption{The state of Erd\H{o}s and Ulam's problem after this paper. Only
the second row of the second conjecture is fully closed; every other row
records either a gap between a proved upper and a proved lower bound, or
an outright open question, and none of the further theorems of this
paper depend on any of the open rows.}
\label{tab:summary}
\end{table}

The one item in the table that is genuinely settled, Erd\H{o}s's second
conjecture, was open, so far as the sources in Section~\ref{sec:related}
record, from 1978 until the argument of Theorem~\ref{thm:F}; the
argument is short (Lemma~\ref{lem:indep} plus a first-moment count) once
it is seen that the two-colouring, forced-monochromatic-union structure
of the problem is exactly what an independent system detects. The first
conjecture is the harder of Erd\H{o}s's two: even the far more permissive
class of cardinality colourings needs the full strength of van der
Waerden's theorem (Theorem~\ref{thm:howorka}) to rule out an $n^c$
bound, and for general colourings no argument in this paper, or, so far
as the authors are aware, in the literature, gives any superlinear lower
bound on $F(n)$ at all beyond the trivial $\lceil(n+1)/2\rceil$. The gap
between what a random colouring is forced to contain
(Theorem~\ref{thm:limits}(ii), exponent $\log_2\log_2n$) and what the
Sauer--Shelah bound permits (Theorem~\ref{thm:F}, exponent $\log_2n$) is
the precise sense in which the first conjecture remains open: closing it
either upward, by a better construction than the random colouring, or
downward, by a bound sharper than Sauer--Shelah, would resolve Question~2
below, and either direction looks, to the authors, to need an idea not
present in this paper.

\section{Questions}
\label{sec:questions}

\begin{question}
Is $f(n)=\lceil(n+1)/2\rceil$ for all $n$? If not, what is the least $n$
for which $f(n)>\lceil(n+1)/2\rceil$, and is $f(n)=O(n)$?
\end{question}

\begin{question}
Is $F(n)$ bounded by a polynomial in $n$? By Theorem~\ref{thm:limits}(ii)
a positive answer cannot come from a random colouring; by
Theorem~\ref{thm:F} the true exponent of $\log n$ in $\log F(n)$ is at
most $2$.
\end{question}

\begin{question}
Is $F(n)>n^{c}$ for every constant $c$ and all large $n$, as Erd\H{o}s
conjectured? The computed values $F(5)=4$, $F(6)=5$, $F(7)=7$ show that
$F$ pulls away from $f$ early, but no lower bound better than
$\lceil(n+1)/2\rceil$ is known for general colourings; for cardinality
colourings Theorem~\ref{thm:howorka} settles the question.
\end{question}

\begin{remark}[An approach that does not work, and precisely why]
\label{rem:posetramsey}
A natural strategy for the first conjecture is to import a Ramsey
theorem for the Boolean lattice itself. The \emph{poset Ramsey number}
$R(Q_k,Q_k)$ is the least $n$ such that every $2$-colouring of the
elements of $2^{[n]}$ contains a monochromatic \emph{copy} of the
Boolean lattice $Q_k=2^{[k]}$, where a copy means an order-embedding
$\phi:2^{[k]}\to2^{[n]}$, that is, an injection with $A\subseteq
B\Leftrightarrow\phi(A)\subseteq\phi(B)$. Axenovich and Walzer
\cite{AW17} proved $R(Q_k,Q_k)\ge2k$, and Lu and Thompson \cite{LT22}
proved $R(Q_k,Q_k)\le k^2-k+2$, so this Ramsey number is polynomial in
$k$. If a monochromatic copy of $Q_k$ were automatically a monochromatic
union-closed family in the sense of this paper, the polynomial bound
would immediately finish the first conjecture: taking $k$ with
$k^2-k+2\le n$, i.e. $k=\Theta(\sqrt n)$, every colouring of $2^{[n]}$
would contain a monochromatic union-closed family of size $2^k=
2^{\Theta(\sqrt n)}$, which exceeds $n^c$ for every fixed $c$ once $n$
is large. This does not work, and the reason is exact rather than a
matter of degree: an order-embedding need not be a lattice embedding.
The constructions in \cite{AW17,LT22} produce $\phi$ by a chain
decomposition of $2^{[n]}$ into symmetric chains, walking up the chains
level by level and, at each level, selecting an available element of the
right colour \emph{from whichever chain happens to offer one}; nothing
in the argument arranges for $\phi(A)\cup\phi(B)$ to equal $\phi(A\cup
B)$ for $A,B\subseteq[k]$, or even for $\phi(A)\cup\phi(B)$ to be
monochromatic, let alone to lie in the image of $\phi$ at all. A copy of
$Q_k$ obtained this way witnesses $2^k$ elements of $2^{[n]}$ correctly
ordered by inclusion and of one colour; it does not witness a
union-closed family, since union-closure is a statement about actual
unions in $2^{[n]}$, not about the abstract order type of a chosen set
of representatives. Repairing this would require either forcing the
chain-selection to respect a fixed block structure throughout, which the
published proofs do not do and which does not obviously survive the
level-by-level, chain-local nature of the argument, or a genuinely
different construction; the authors were not able to find one, and are
not aware of one in the literature. The gap between the two notions of
``copy'' is therefore, at present, exactly the gap between what is known
about poset Ramsey numbers and what the first conjecture requires.
\end{remark}

\begin{remark}[A second approach that does not work, and precisely why]
\label{rem:booleanalgebra}
Unlike the previous approach, there is a construction that genuinely
avoids the defect of Remark~\ref{rem:posetramsey}: pairwise disjoint
nonempty $X_1,\dots,X_d\subseteq[n]$ generate, by
$I\mapsto\bigcup_{i\in I}X_i$, an actual sublattice of $2^{[n]}$ isomorphic
to $2^{[d]}$, a \emph{$d$-dimensional Boolean algebra} in the terminology
of Gunderson, R\"odl and Sidorenko \cite{GRS99}; if this copy is
monochromatic, its $2^d-1$ nonempty members are a genuine monochromatic
union-closed family (indeed a sublattice), with no gap between ``copy''
and ``union-closed family'' to repair. The relevant Ramsey function,
$r(d,n)$, is the largest $r$ such that every $r$-colouring of the subsets
of an $n$-set contains a monochromatic $d$-dimensional Boolean algebra;
Gunderson, R\"odl and Sidorenko prove
\[
c_1\,n^{1/2^d}\ \le\ r(d,n)\ \le\ n^{\,d/2^{d-1}(1+o(1))}
\]
for an absolute constant $c_1>0$ (\cite[Theorem 4.7]{GRS99}; the case
$d=2$ is their Theorem 4.6, with the sharper constants
$(1-o(1))\tfrac34\sqrt n\le r(2,n)\le(1+o(1))\sqrt n$). Fixing $r=2$ and
asking how large a $d$ this forces at a given $n$ is exactly the question
needed for the first conjecture, and both bounds answer it the same way:
the lower bound shows $r(d,n)\ge2$, so that the guarantee holds, as soon
as $n\ge(2/c_1)^{2^d}$, while the upper bound shows a $2$-colouring
avoiding every monochromatic $d$-dimensional Boolean algebra exists as
soon as $n^{2d/2^d(1+o(1))}<2$, i.e.\ for $n$ below essentially the same
threshold. Both directions therefore agree that the threshold is
\emph{doubly} exponential in $d$, $n=\Theta(2^{2^d})$ up to the constant in
the exponent, not tower-type as in Remark~\ref{rem:posetramsey}'s more
distant relative, the Hilbert cube problem \cite{Gow01,Bal19}, but also
not singly exponential. Inverting, the largest $d$ that a $2$-colouring
of $2^{[n]}$ is guaranteed to force by this method is only
$d(n)=(1+o(1))\log_2\log_2n$, and the resulting monochromatic
union-closed family has only $2^{d(n)}-1=\Theta(\log n)$ members: smaller,
for every large $n$, than the trivial chain bound
$\lceil(n+1)/2\rceil$ of \eqref{eq:trivial}. Because the upper bound in
\cite{GRS99} matches the lower bound in the same doubly-exponential
shape, this is not a matter of a Ramsey number not yet being computed
well enough: for a fixed pair of colours, monochromatic Boolean
subalgebras generated by disjoint blocks are, provably, too weak a
structure to beat a single chain, let alone to grow polynomially with
$n$. Any resolution of the first conjecture must therefore use
union-closed families that are not sublattices of this particularly
symmetric, block-generated kind.
\end{remark}

Both remarks above attack the first conjecture directly, over all of
$2^{[n]}$ at once. The last attempt instead tries to build it up by
recursion on $n$, and the recursion turns out to be exactly as hard as
the problem it is meant to simplify -- provably so, not just in practice.

\begin{lemma}[Doubling, conditional on a simultaneous family]
\label{lem:doubling}
Let $n\ge1$, let $\chi$ be a $2$-colouring of $2^{[n]}$, and define
$\chi_0,\chi_1:2^{[n-1]}\to\{0,1\}$ by $\chi_0(A)=\chi(A)$ and
$\chi_1(A)=\chi(A\cup\{n\})$. If there is a union-closed family
$F\subseteq2^{[n-1]}$ and a colour $c$ such that every member of $F$ has
colour $c$ under \emph{both} $\chi_0$ and $\chi_1$, then
\[
G\ =\ F\ \cup\ \{A\cup\{n\}:A\in F\}
\]
is a monochromatic (colour $c$) union-closed family for $\chi$, and
$|G|=2|F|$.
\end{lemma}

\begin{proof}
$G$ has colour $c$ throughout, by hypothesis. For closure under union,
take $P,Q\in G$. If $P,Q\in F$ then $P\cup Q\in F\subseteq G$ since $F$
is union-closed. If $P=A\cup\{n\}$ and $Q=A'\cup\{n\}$ with $A,A'\in F$,
then $P\cup Q=(A\cup A')\cup\{n\}\in G$ since $A\cup A'\in F$. If
$P=B\in F$ and $Q=A\cup\{n\}$ with $A\in F$, then
$P\cup Q=(A\cup B)\cup\{n\}\in G$ since $A\cup B\in F$. The two halves of
$G$ are disjoint, since one avoids $n$ and the other contains it, so
$|G|=2|F|$.
\end{proof}

If the lemma's hypothesis -- a single union-closed family simultaneously
monochromatic, in the same colour, under two given colourings of
$2^{[n-1]}$ -- could always be met with $|F|\ge F(n-1)$, then applying it
at every level from some fixed $n_0$ up would give
$F(n)\ge2^{\,n-n_0}F(n_0)$ for all $n\ge n_0$, which is $2^{\Omega(n)}$.
This contradicts Theorem~\ref{thm:F}: $F(n)\le n^{(1+o(1))\log_2n}
=2^{(1+o(1))(\log_2n)^2}=2^{o(n)}$. Hence the required simultaneous
family fails to exist, with $|F|\ge F(n-1)$, for infinitely many $n$ -- if
it failed only finitely often, the recursion could still be run from the
last failure onward and the same contradiction would follow. This is not
an empirical difficulty with the three patches attempted below; it is a
proof, from the paper's own Theorem~\ref{thm:F}, that no uniform way of
supplying the lemma's hypothesis at every level can exist.

\begin{remark}[A third approach that does not work, and precisely why]
\label{rem:doubling}
The obstruction identified above still leaves open whether the
simultaneous family can be found \emph{often enough} -- say, at all $n$
outside a sparse exceptional set -- to salvage a weaker recursion; the
authors tried three routes to supplying it and each one breaks
union-closure instead of repairing it. First, apply the first conjecture's
guarantee to $\chi_0$ alone to get a union-closed $F_0\subseteq2^{[n-1]}$
of size $\ge F(n-1)$ in some colour $c_0$, then $2$-colour $F_0$ itself by
$\chi_1$: by pigeonhole at least $|F_0|/2$ members of $F_0$ share a single
$\chi_1$-colour $c_1$, but this majority subfamily need not be
union-closed, since the union of two of its members can fall on the
$\chi_1$-minority side. Second, treat $\psi=(\chi_0,\chi_1)$ as a single
$4$-colouring of $2^{[n-1]}$ and ask for a union-closed family
monochromatic under $\psi$: this only helps if the forced $\psi$-colour
happens to be a diagonal value $(c,c)$ rather than $(0,1)$ or $(1,0)$, and
nothing forces the diagonal. Third, split $[n]$ into two halves instead
of peeling one point, and look slice by slice: each slice $\chi_A$
(for $A$ ranging over subsets of the first half) admits its own
union-closed family by induction on half the size, but combining these
into one union-closed family over all of $[n]$ requires many
independently-obtained structures to agree, which is the same difficulty
that gives Remark~\ref{rem:booleanalgebra} its doubly-exponential cost, not
a way around it. The recursion of Lemma~\ref{lem:doubling} is therefore
not merely unproved at every step; by the argument above it is
\emph{provably} unavailable at infinitely many steps, and the three
natural repairs available to a coloring-agnostic argument all either
destroy closure or reduce to the difficulty already isolated in
Remark~\ref{rem:booleanalgebra}.
\end{remark}

\section{Formal verification in Lean}
\label{sec:lean}

The definitions of $f(n)$ and $F(n)$, the first assertion of
Theorem~\ref{thm:F}, the resulting resolution of Erd\H{o}s's second
conjecture, and the trivial lower bound \eqref{eq:trivial} have been
formalised in the Lean~4 proof assistant \cite{MU21} on top of the
Mathlib library \cite{Mathlib20}. The formalisation is the directory
\texttt{lean/} of the repository accompanying the paper
(\url{https://github.com/creelie/erdos-1183}); it builds with Lean~4.34.1
and Mathlib~v4.34.1, contains no \texttt{sorry} and declares no axiom, and
for every theorem listed below Lean's \texttt{\#print axioms} reports only
the three standard axioms \texttt{propext}, \texttt{Classical.choice} and
\texttt{Quot.sound}.

In the formalisation $[n]$ is \texttt{Fin n}, a colouring is a function
\texttt{Finset (Fin n) → Bool}, and \texttt{bigF n} and
\texttt{smallF n} are defined literally as in Section~\ref{sec:intro}, as
the minimum over colourings of the largest monochromatic family closed
under union (respectively under union and intersection). The lemmas
\texttt{le\_bigF\_iff} and \texttt{bigF\_le\_iff} check that
\texttt{bigF n} is exactly the largest $m$ such that every colouring
contains a monochromatic union-closed family of $m$ sets.

\begin{table}[h]
\centering
\begin{tabular}{p{0.34\linewidth} p{0.56\linewidth}}
\toprule
Lean theorem & Statement\\
\midrule
\texttt{bigF\_le\_sum\_choose} & if $nd+2<2^d$ then
$F(n)\le\sum_{j<d}\binom nj$ (Theorem~\ref{thm:F}, first assertion)\\
\addlinespace
\texttt{bigF\_le\_explicit} & for $n\ge2$ and $d=2(\lfloor\log_2n\rfloor+1)$,
$F(n)\le d\,n^{d}$\\
\addlinespace
\texttt{second\_conjecture} & for every real $\eps>0$,
$F(n)<(1+\eps)^n$ for all sufficiently large $n$\\
\addlinespace
\texttt{half\_le\_smallF}, \texttt{smallF\_le\_bigF} &
$\lceil(n+1)/2\rceil\le f(n)\le F(n)$, i.e.\ \eqref{eq:trivial}\\
\bottomrule
\end{tabular}
\caption{Machine-checked statements.}
\label{tab:lean}
\end{table}

The formal proof follows Section~\ref{sec:F} with two changes. The
first-moment argument is replaced by the equivalent count: for a fixed
independent $d$-tuple the $2^d-1$ unions are distinct
(Lemma~\ref{lem:indep}(a)), so at most $2\cdot2^{2^n-(2^d-1)}$ colourings
make the tuple bad, and summing over the at most $2^{nd}$ tuples gives
fewer than $2^{2^n}$ colourings when $nd+2<2^d$. The Sauer--Shelah lemma
is Mathlib's (\texttt{Finset.card\_le\_card\_shatterer} and
\texttt{Finset.card\_shatterer\_le\_sum\_vcDim}). For the asymptotic
statement the formalisation uses the simpler choice
$d=2(\lfloor\log_2n\rfloor+1)$, which gives $F(n)\le n^{O(\log n)}$ and
suffices for $F(n)<(1+\eps)^n$; the sharper exponent
$(1+o(1))\log_2n$ of Theorem~\ref{thm:F} is not formalised.

Theorems~\ref{thm:f}, \ref{thm:limits}, \ref{thm:small} and
\ref{thm:howorka} are not formalised; Theorem~\ref{thm:small} rests on the
satisfiability computations of Section~\ref{sec:small}. Erd\H{o}s's first
conjecture is stated in the formalisation as the proposition
\texttt{FirstConjecture} ($F(n)>n^{c}$ for every fixed $c$ and all large
$n$) and is not proved: it remains open.

\appendix
\section{The satisfiability encoding}
\label{app:sat}

This appendix formalises the algorithm sketched in
Section~\ref{sec:small}, for a reader who wants to reproduce or
re-implement it without reading the code.

\subsection{Variables and the colouring clauses}
Fix $n$. The Boolean variables are $x_S$ for $S\subseteq[n]$, one per
subset, with $x_S=1$ meaning ``$S$ has colour $1$''. The symmetry
$\chi\mapsto1-\chi$ is broken once and for all by the unit clause
$\lnot x_\emptyset$. No other clause is imposed at the outset; clauses
are added only as counterexamples are found, in the loop described next.

\subsection{The closed-family clause}
For a family $\cF\subseteq2^{[n]}$ with $|\cF|=M\ge2$, forbidding $\cF$
from being monochromatic is the pair of clauses
\[
\Bigl(\bigvee_{S\in\cF}\lnot x_S\Bigr)\qquad\text{and}\qquad
\Bigl(\bigvee_{S\in\cF}x_S\Bigr),
\]
the first true unless every $S\in\cF$ has $x_S=1$, the second true
unless every $S\in\cF$ has $x_S=0$; together they say that $\cF$ is not
entirely one colour. There are $2^{2^n}-2$ closed families of size
$\ge2$ in the worst case (Lemma~\ref{lem:count} bounds this far more
sharply), far too many to write down in advance for $n\ge5$ or so; the
counterexample-guided loop adds only the ones that are actually needed.

\subsection{The outer loop}
For a target $M$, the loop maintains a growing CNF formula $\Phi$,
initialised to $\{\lnot x_\emptyset\}$.
\begin{enumerate}[label=\textup{(\arabic*)}]
\item Call the SAT solver on $\Phi$. If unsatisfiable, report that every
colouring has a monochromatic closed family of size $\ge M$ (so the
value being computed is $\ge M$) and stop.
\item Otherwise let $\chi$ be the returned model. For each colour
$c\in\{0,1\}$, run the inner search of Section~\ref{sec:inner} below on
the colour class $\chi^{-1}(c)$ with target $M$.
\item If neither inner search finds a closed family of size $\ge M$,
$\chi$ is a witness that the value is $<M$; report $\chi$ and stop.
\item Otherwise, for each colour class in which the inner search found a
closed family $\cF$ of size $\ge M$, add the two clauses of
Section~\ref{sec:sub-clause} for that $\cF$ (or, if $|\cF|>M$, for an
$M$-element closed subfamily of it, which exists since removing a
maximal element of a closed family under either operation leaves a
smaller closed family) to $\Phi$, and go to (1).
\end{enumerate}
Termination follows because each pass through the loop either stops or
adds at least one clause pair ruling out at least one previously
satisfying model, and there are only $2^{2^n}$ models in total; in
practice the loop terminates after examining a number of families in
the tens or low hundreds of thousands (Section~\ref{sec:small} records
the exact counts for the two largest cases, $f(9)$ and $F(7)$), far
fewer than $2^{2^n}$.

\subsection{The inner search}
\label{sec:inner}
Given a set $C\subseteq2^{[n]}$ (a colour class) and a target $M$, the
inner search looks for a subfamily $\cF\subseteq C$ with $|\cF|\ge M$
that is closed under the relevant operations (union alone for $F$,
union and intersection for $f$). It performs a depth-first
branch-and-bound search that maintains a partial closed family, at each
step either adding an unused element of $C$ and closing (repeatedly
adjoining unions, and for $f$ also intersections, of existing members
that lie in $C$, until no new member of $C$ is produced) or backtracking;
a branch is pruned as soon as the sum of the current family's size and
the number of remaining candidates in $C$ falls below $M$. This is
exactly the sense in which $\cF$ is ``built by adding sets one at a
time and closing under the operations'' in Section~\ref{sec:small}.

\subsection{The independent verifier}
\label{sec:sub-clause}
The witness-checking program of Section~\ref{sec:small} re-expresses,
for a fixed colour class $C$ and target $M$, the existence of a closed
subfamily of size $\ge M$ as a second, differently structured
satisfiability instance: one Boolean variable $y_S$ for each $S\in C$,
the closure conditions $\lnot y_S\lor\lnot y_{S'}\lor y_{S\cup S'}$ (and,
for $f$, the analogous clause for $S\cap S'$) for all $S,S'\in C$, and a
cardinality constraint $\sum_{S\in C}y_S\ge M$ encoded by the sequential
counter method of Sinz \cite{Si05}, which introduces $O(|C|M)$
auxiliary variables and clauses rather than the $\binom{|C|}M$ clauses of
a naive encoding. When $S\cup S'\notin C$ (respectively $S\cap S'\notin
C$), the corresponding closure clause degenerates to $\lnot
y_S\lor\lnot y_{S'}$, since no third variable is available to witness
closure and $S,S'$ simply cannot be chosen together; this is the same
convention used, without comment, in the code. This instance shares no
variables, clauses, or code with the branch-and-bound search of
Section~\ref{sec:inner}; agreement
between the two, checked for every entry of Theorem~\ref{thm:small}
(Section~\ref{sec:small}), is what makes the values independently
verified rather than merely computed once.

\section*{Competing Interests}
The authors declare that there are no competing interests, financial or
non-financial, that could have influenced the research reported in this
paper.

\section*{Data Availability}
The program computing the values in Theorem~\ref{thm:small}, with its
output, and the Lean formalisation of Section~\ref{sec:lean} accompany the
paper and are available at \url{https://github.com/creelie/erdos-1183}. The underlying data and code are also
available from the corresponding authors upon reasonable request by
email.

\begin{thebibliography}{AHS24}
\raggedright

\bibitem[AHS24]{AHS24}
R.~Alweiss, B.~Huang, and M.~Sellke,
\emph{Improved lower bound for Frankl's union-closed sets conjecture},
Electron. J. Combin. \textbf{31} (2024), no.~3, \#P3.35.
doi:10.37236/12232

\bibitem[AW17]{AW17}
M.~Axenovich and A.~Walzer,
\emph{Boolean lattices: Ramsey properties and embeddings},
Order, published online 2016/2017, arXiv:1512.05565.
doi:10.1007/s11083-016-9399-7

\bibitem[Bl26]{Bl26}
T.~F. Bloom,
\emph{Erd\H{o}s Problem \#1183},
\url{https://www.erdosproblems.com/1183}, accessed 6 September 2026.

\bibitem[Bo86]{Bo86}
B.~Bollob\'as,
\emph{Combinatorics: Set Systems, Hypergraphs, Families of Vectors and
Combinatorial Probability},
Cambridge University Press, Cambridge, 1986.

\bibitem[Di50]{Di50}
R.~P. Dilworth,
\emph{A decomposition theorem for partially ordered sets},
Ann. of Math. (2) \textbf{51} (1950), 161--166.
doi:10.2307/1969503

\bibitem[Er78]{Er78}
P.~Erd\H{o}s,
\emph{Problems and results in combinatorial analysis and combinatorial
number theory},
in: Proceedings of the Ninth Southeastern Conference on Combinatorics,
Graph Theory, and Computing (Boca Raton, 1978), Congressus Numerantium
XXI, Utilitas Math., Winnipeg, 1978, pp.~29--40.

\bibitem[Fr79]{Fr79}
P.~Frankl's conjecture, first published in D.~Duffus,
\emph{Open problem session}, in: Graphs and Order (I.~Rival, ed.), NATO
Advanced Science Institutes Series C, vol.~147, D.~Reidel, Dordrecht,
1985, p.~525.

\bibitem[Bal19]{Bal19}
J.~Balogh, M.~Lavrov, G.~Shakan, and A.~Zs.~Wagner,
\emph{Monochromatic Hilbert cubes and arithmetic progressions},
Electron. J. Combin. \textbf{26} (2019), no.~2, \#P2.22.
doi:10.37236/7917

\bibitem[Gi22]{Gi22}
J.~Gilmer,
\emph{A constant lower bound for the union-closed sets conjecture},
arXiv:2211.09055, 2022.

\bibitem[Gow01]{Gow01}
W.~T.~Gowers,
\emph{A new proof of Szemer\'edi's theorem},
Geom. Funct. Anal. \textbf{11} (2001), 465--588.
doi:10.1007/s00039-001-0332-9

\bibitem[GRS99]{GRS99}
D.~S.~Gunderson, V.~R\"odl, and A.~Sidorenko,
\emph{Extremal problems for sets forming Boolean algebras and complete
partite hypergraphs},
J. Combin. Theory Ser. A \textbf{88} (1999), 342--367.
doi:10.1006/jcta.1999.2973

\bibitem[LT22]{LT22}
L.~Lu and J.~C.~Thompson,
\emph{Poset Ramsey numbers for Boolean lattices},
Order \textbf{39} (2022), 171--185.
doi:10.1007/s11083-021-09557-4

\bibitem[IMM18]{IMM18}
A.~Ignatiev, A.~Morgado, and J.~Marques-Silva,
\emph{PySAT: a Python toolkit for prototyping with SAT oracles},
in: Theory and Applications of Satisfiability Testing (SAT 2018),
Lecture Notes in Computer Science 10929, Springer, 2018, pp.~428--437.
doi:10.1007/978-3-319-94144-8\_26

\bibitem[MU21]{MU21}
L.~de~Moura and S.~Ullrich,
\emph{The Lean 4 theorem prover and programming language},
in: Automated Deduction -- CADE 28, Lecture Notes in Computer Science
12699, Springer, 2021, pp.~625--635.
doi:10.1007/978-3-030-79876-5\_37

\bibitem[Mat20]{Mathlib20}
The mathlib Community,
\emph{The Lean mathematical library},
in: Proceedings of the 9th ACM SIGPLAN International Conference on
Certified Programs and Proofs (CPP 2020), ACM, 2020, pp.~367--381.
doi:10.1145/3372885.3373824

\bibitem[Sa72]{Sa72}
N.~Sauer,
\emph{On the density of families of sets},
J. Combin. Theory Ser. A \textbf{13} (1972), 145--147.
doi:10.1016/0097-3165(72)90019-2

\bibitem[Sh72]{Sh72}
S.~Shelah,
\emph{A combinatorial problem; stability and order for models and
theories in infinitary languages},
Pacific J. Math. \textbf{41} (1972), 247--261.
doi:10.2140/pjm.1972.41.247

\bibitem[Si05]{Si05}
C.~Sinz,
\emph{Towards an optimal CNF encoding of Boolean cardinality
constraints},
in: Principles and Practice of Constraint Programming (CP 2005),
Lecture Notes in Computer Science 3709, Springer, 2005, pp.~827--831.
doi:10.1007/11564751\_73

\bibitem[VC71]{VC71}
V.~N. Vapnik and A.~Ya. Chervonenkis,
\emph{On the uniform convergence of relative frequencies of events to
their probabilities},
Theory Probab. Appl. \textbf{16} (1971), 264--280.
doi:10.1137/1116025

\bibitem[vdW27]{vdW27}
B.~L. van der Waerden,
\emph{Beweis einer Baudetschen Vermutung},
Nieuw Arch. Wiskd. \textbf{15} (1927), 212--216.

\bibitem[St12]{St12}
R.~P. Stanley,
\emph{Enumerative Combinatorics, Volume 1},
second edition, Cambridge Studies in Advanced Mathematics 49, Cambridge
University Press, Cambridge, 2012.
doi:10.1017/CBO9781139058520

\end{thebibliography}

\end{document}
